Bộ dữ liệu California Housing được xây dựng từ cuộc điều tra dân số Hoa Kỳ năm
1990, trong đó mỗi quan sát tương ứng với một nhóm khu vực điều tra dân số
(census block group), là đơn vị địa lý nhỏ nhất trong thống kê dân số, thường
bao gồm từ vài trăm đến vài nghìn người \cite{pace1997sparse}.

Dữ liệu được thu thập trên toàn bang California với 20,640 quan sát và 9 đặc
trưng. Mỗi nhóm khu vực được xem như một vùng dân cư và được đại diện bởi một
điểm trung tâm (centroid) nhằm phục vụ cho phân tích không gian. Khoảng cách
giữa các khu vực được tính dựa trên tọa độ địa lý (kinh độ và vĩ độ) của các
điểm trung tâm này, từ đó phản ánh mức độ liên hệ không gian giữa các vùng
\cite{pace1997sparse}.

Trung bình, mỗi nhóm khu vực có khoảng 1,425 cư dân sinh sống trong một khu vực
địa lý nhỏ. Do đặc điểm phân bố dân cư, diện tích các khu vực thường giảm khi
mật độ dân số tăng \cite{pace1997sparse}.

Trong quá trình xử lý dữ liệu, các quan sát có giá trị bằng 0 ở cả biến đầu vào
và biến mục tiêu đã được loại bỏ để đảm bảo tính nhất quán
\cite{pace1997sparse}.

Biến mục tiêu là giá trị trung vị nhà ở (median house value) trong từng khu
vực, được đo theo đơn vị 100,000 USD. Một số đặc trưng như số phòng và số phòng
ngủ được tính trung bình theo hộ gia đình (household), vì vậy có thể xuất hiện
giá trị lớn ở những khu vực có ít hộ dân nhưng nhiều nhà trống, chẳng hạn khu
nghỉ dưỡng.
